\chapter{Verilog 구문을 실제 회로로 번역하는 법}

\begin{summarybox}[이 장의 관점]
RTL을 읽을 때에는 문법의 뜻과 생성되는 회로를 동시에 말할 수 있어야 한다.
\rtlfile{assign}은 조합 회로, clocked \rtlfile{always}의 nonblocking assignment는
register, 조건문은 data 또는 enable을 고르는 MUX가 된다. loop는 시간을 두고
반복 실행되는 문장이 아니라 elaboration 과정에서 펼쳐지는 병렬 하드웨어다.
\end{summarybox}

\section{Continuous assignment와 combinational cone}

\begin{lstlisting}
wire [22:0] x = (a ^ b) + c;
\end{lstlisting}

이 코드는 XOR를 한 다음 다음 cycle에 add하는 순차 동작이 아니다. \rtlfile{a},
\rtlfile{b}, \rtlfile{c}가 바뀌면 XOR와 carry chain을 거쳐 \rtlfile{x}가 안정되는
하나의 combinational cone이다. 중간 wire를 추가해도 register가 아니므로 timing
경계는 생기지 않는다.

\begin{lstlisting}
wire [22:0] xor_w = a ^ b;
wire [22:0] x     = xor_w + c;
\end{lstlisting}

두 번째 표현은 debug와 report tracing에는 도움이 되지만 첫 번째 표현과 같은
회로로 최적화될 수 있다. pipeline을 만들려면 clock edge에서 값을 저장해야 한다.

\section{Clocked always block과 register input MUX}

\begin{lstlisting}
always @(posedge clk) begin
  if (rst)
    q <= 23'd0;
  else if (en)
    q <= d;
end
\end{lstlisting}

이 블록은 \rtlfile{q}의 모든 bit에 FF를 만든다. \rtlfile{rst}와 \rtlfile{en}은
코드가 위에 있다고 먼저 실행되는 연산이 아니라 FF의 reset/enable 기능 또는
D input 앞 MUX로 구현된다. exact mapping은 reset 형태와 target primitive,
optimization에 따라 달라진다.

clocked block 안에서 여러 nonblocking assignment가 있으면 오른쪽 식은 모두
현재 cycle의 값을 읽고, 왼쪽 register는 같은 edge 뒤에 함께 갱신된다.

\begin{lstlisting}
always @(posedge clk) begin
  a_q <= in;
  b_q <= a_q;
end
\end{lstlisting}

여기서 \rtlfile{b\_q}는 새 \rtlfile{in}이 아니라 이전 \rtlfile{a\_q}를 받는다.
따라서 두 stage pipeline이 된다. 이 cycle 의미를 이해해야 data와 valid, mode,
address를 같은 stage로 이동시킬 수 있다.

\section{Blocking과 nonblocking: 회로 종류보다 stage 의미}

조합 block에서는 blocking assignment로 위에서 계산한 임시 값을 아래 식에
사용하는 것이 자연스럽다.

\begin{lstlisting}
always @* begin
  sum_w = a + b;
  if (sum_w >= Q)
    out_w = sum_w - Q;
  else
    out_w = sum_w;
end
\end{lstlisting}

이 코드는 add와 correction이 같은 cycle에 직렬로 놓인다는 점이 핵심이다.
blocking assignment를 썼기 때문에 software처럼 시간이 흐르는 것이 아니다.

clocked block에는 nonblocking assignment를 사용해 각 register stage의 이전 값과
다음 값을 분명히 한다.

\begin{lstlisting}
always @(posedge clk) begin
  raw_q <= a + b;
  if (raw_q >= Q)
    corr_q <= raw_q - Q;
  else
    corr_q <= raw_q;
end
\end{lstlisting}

이 표현은 add와 correction을 서로 다른 register 경계에 둔다. latency가 한
cycle 늘고 throughput은 매 cycle 새 입력을 받을 수 있다. control signal도 같은
cycle만큼 지연해야 한다.

\section{if, case와 MUX 구조}

조건에 따라 하나의 register 또는 wire에 여러 값을 대입하면 data selection이
필요하다.

\begin{lstlisting}
always @* begin
  case (mode)
    2'd0: y = a;
    2'd1: y = b;
    2'd2: y = c;
    default: y = d;
  endcase
end
\end{lstlisting}

배타적인 네 선택지는 일반적으로 4:1 MUX 구조다. 반면 다음 코드는 첫 조건이
가장 높은 우선순위를 가져야 하므로 priority chain이 된다.

\begin{lstlisting}
always @* begin
  if      (req[0]) y = d0;
  else if (req[1]) y = d1;
  else if (req[2]) y = d2;
  else             y = d3;
end
\end{lstlisting}

priority가 기능 요구라면 이 구조가 맞다. 선택지가 원래 one-hot 또는 mutually
exclusive라면 그 계약을 assertion으로 검증하고 평행 selection이 드러나는
구조를 쓸 수 있다. 단순히 \rtlfile{case}가 언제나 빠르다고 외우지 않고 기능적
우선순위가 필요한지 먼저 판단한다.

\section{조합 block의 incomplete assignment와 latch}

\begin{lstlisting}
always @* begin
  if (en)
    y = d;
end
\end{lstlisting}

\rtlfile{en=0}일 때 \rtlfile{y} 값을 정의하지 않았으므로 이전 값을 유지해야 한다.
clock이 없는데 상태를 보존하려면 level-sensitive latch가 필요하다. 의도하지 않은
latch는 timing 분석과 검증을 복잡하게 만든다.

\begin{lstlisting}
always @* begin
  y = 23'd0;
  if (en)
    y = d;
end
\end{lstlisting}

기본값을 먼저 주면 모든 입력 조건에서 값이 정의되고 MUX로 합성된다. latch가
정말 필요한 회로라면 의도를 명시하고, 대부분의 synchronous datapath에서는
조합 output을 완전히 대입한다.

\section{for loop와 generate: 공간 복제인지 시간 반복인지}

\begin{lstlisting}
integer i;
always @* begin
  for (i = 0; i < 4; i = i + 1)
    y[i] = a[i] ^ b[i];
end
\end{lstlisting}

고정 횟수 loop는 네 번의 XOR 하드웨어로 펼쳐진다. 한 개의 XOR가 네 cycle 동안
재사용되는 구조가 아니다. $N$개의 butterfly를 generate하면 대체로 $N$개의
datapath가 생긴다. 시간 공유를 원하면 counter/FSM, operand MUX, result register와
latency protocol을 직접 설계해야 한다.

\begin{summarybox}[병렬도와 면적의 첫 근사]
폭 $W$의 operator를 $P$개 병렬 생성하면 일반 fabric 비용은 대략
\[
  A \propto W\times P\times C_{op}
\]
에서 시작한다. 실제 합성에서는 constant propagation, resource sharing, DSP/BRAM
mapping과 packing이 이 값을 바꾼다. loop bound와 generate count는 면적을 읽을 때
가장 먼저 확인할 숫자다.
\end{summarybox}

\section{Reduction과 bitwise/logical operator}

다음 세 표현은 다른 회로다.

\begin{lstlisting}
wire [22:0] bit_or = a | b;
wire        red_or = |a;
wire        log_or = a || b;
\end{lstlisting}

\rtlfile{a|b}는 23개의 1-bit OR가 병렬로 존재한다. \rtlfile{|a}는 23 bit를
하나로 모으는 reduction tree다. \rtlfile{a||b}는 두 bus가 각각 zero인지 판단한
뒤 두 Boolean 결과를 OR한다. reduction 결과가 넓은 enable과 memory selection을
구동하면 작은 LUT count와 별개로 fanout/route 병목이 될 수 있다.

Unified NTT에서 13-stage valid vector의 reduction OR는 isolated cost가 3 LUT에
불과했다. 하지만 그 결과가 BUSY와 memory phase control까지 연결되어 있었기
때문에 occupancy counter로 상태를 local하게 만든 것이 timing에 의미가 있었다.

\section{비트 폭: 저장 면적과 arithmetic 깊이를 동시에 결정한다}

Verilog expression의 폭을 암묵적으로 두면 의도하지 않은 truncation 또는
extension이 생길 수 있다.

\begin{lstlisting}
wire [22:0] a, b;
wire [23:0] full_sum = {1'b0, a} + {1'b0, b};
wire [22:0] wrap_sum = a + b;
\end{lstlisting}

첫 식은 carry-out까지 보존한다. 두 번째는 23 bit에서 wrap된다. 모듈러 연산에서
어느 범위를 보존해야 하는지 증명한 뒤 최소 폭을 선택한다. 한 bit를 줄이면 해당
register, MUX, adder, comparator와 배선에서 반복적으로 자원이 줄 수 있다.

반대로 필요한 guard bit를 제거하면 일부 입력에서만 틀리는 오류가 생긴다. 폭
최적화는 simulation 몇 개가 아니라 operand range와 intermediate bound를 근거로
해야 한다.

\section{signed와 unsigned: 같은 bit pattern을 다르게 해석한다}

FPGA 배선에는 부호 표지가 없다. signedness는 comparison, extension, shift와
arithmetic 식을 합성하는 규칙이다. \rtlfile{4'b1111}은 unsigned 15이기도 하고
signed $-1$이기도 하다.

\begin{lstlisting}
wire signed [8:0] diff =
    $signed({1'b0, a8}) - $signed({1'b0, b8});
wire [7:0] corrected = diff[8] ? diff + Q : diff[7:0];
\end{lstlisting}

unsigned operand 두 개를 빼더라도 음수 여부가 필요하면 한 bit 확장한 signed
intermediate 또는 borrow를 보존해야 한다. signed와 unsigned를 섞은 comparison은
먼저 같은 폭과 의도한 부호로 명시적으로 cast한다.

\begin{warningbox}[상수의 폭도 회로의 일부다]
unsized literal과 signed literal은 expression 폭과 부호 확장에 영향을 줄 수
있다. 특히 parameterized RTL에서는 \rtlfile{23'd...}, \rtlfile{\$signed(...)}와
명시적 concatenation으로 intermediate 폭을 읽을 수 있게 만든다. synthesis
warning이 없다는 사실은 수학적 범위가 맞다는 증명이 아니다.
\end{warningbox}

\section{연산자를 보자마자 그릴 수 있어야 하는 회로}

\begin{longtable}{L{35mm}L{52mm}L{72mm}}
\toprule
RTL 표현 & 우선 떠올릴 회로 & 다음 질문 \\
\midrule
\endhead
\rtlfile{a \& b}, \rtlfile{a \^ b} & bit별 LUT 병렬 구조 & output 폭과 LUT packing은? \\
\rtlfile{|a} & reduction tree & level과 결과 fanout은? \\
\rtlfile{a+b}, \rtlfile{a-b} & LUT + carry chain & 폭과 앞뒤 MUX는? \\
\rtlfile{a==b} & XNOR/reduction 성격 & constant bit가 제거되는가? \\
\rtlfile{a<b} & magnitude/borrow 구조 & signedness와 폭은? \\
ternary, case & bit별 MUX & input 수, priority, select fanout은? \\
\rtlfile{a<<3} & 배선 재배치 & 잘려 나가는 bit와 width는? \\
\rtlfile{a<<shamt} & barrel MUX & shift 후보 수와 stage는? \\
\rtlfile{a*b} & DSP 또는 fabric multiplier & operand 폭, signedness, DSP budget은? \\
clocked assignment & FF 또는 hard-block register & reset/enable과 stage alignment는? \\
fixed loop/generate & 병렬 복제 hardware & unroll 개수와 공유 의도는? \\
array read/write & FF/LUTRAM/BRAM 후보 & port와 synchronous-read template은? \\
\bottomrule
\end{longtable}

이 번역표는 정확한 utilization 표가 아니다. RTL review에서 회로 후보를 빠르게
그린 뒤 synthesis 결과로 수정하기 위한 첫 모델이다.
